% \section{Introduction}
% It is not a funny thing for AI to let AI model understand the joke of funny video. Consider the humorous video in Figure \ref{figure_main}. The boy was looking at his phone normally, but then was accidentally squeezed all over his face with ketchup. The contrast before and after is hilarious. Human can easily realize what is unusual about this video, while to make AI realize that there is something counter-intuitive in this video and explain humor is a challenging task. It needs a model to understand that squeeze ketchup on the face is uncommon and unexpected.\\
% Besides humorous videos, there are many other kinds of videos that need complex reasoning ability to accurately understand the counter-intuitive features of videos. For example, in Japan's Kasou Taishou [quote], people disguise themselves to achieve various unexpected effects, which are highly creative. The surprising effect of magic videos [quote] is also due to its counter-intuitive nature. AI needs to understand where these videos are counter-factual to explain why they are amazing. In other words, if the model can locate, describe and explain the counter-intuitive points on these videos, this will be strong support for the model's reasoning performance.\\
% In the task of picture reasoning, the model usually completes spatial relationship reasoning and infers possible behaviors [] according to experience. Because the complexity of humor comes from the difference between reality and expectation [] as long as the juxtaposition and transformation of incongruous things [], our task needs a model to better understand both temporal and spatial information. In this benchmark, the reasoning for video requires the model to have the ability of Spatial-Temporal reasoning, that is, these task cannot be completed through the pictures of a frame in the video.\\
% In some classic VideoQA datasets, such as TVQA \cite{lei2018tvqa} and MovieQA \cite{mun2017marioqa}, the input includes not only video, but also audio and even subtitle and bounding box. For reasoning on these datasets, the understanding of audio is indispensable. Different from this, we hope to focus on the visual understanding ability of the model. Therefore, our video will not eliminate sound, but the reasoning of our video can be completed only by visual content. More datasets, such as Social-IQ [], also want to perform audio-independence reasoning, but we have deeper requirements for visual understanding of video.\\
% In this paper, we have collected 10k humorous videos with an average duration of 20s, 3k creative super-variety videos and 1k magic videos, and our data set also includes 40k annotation. These counter-intuitive videos challenge AI's reasoning ability.\\
% Next, we propose three main reasoning tasks on the FunQA dataset:
% (1) time stamp localization, where model needs to find out the time stamps of the beginning and end of counter-intuitive segments in video.
% (2) description generation, where model objectively describes the counter-intuitive fragments in the video.
% (3) explanation generation, where the model need to generate detailed explanations on why the video is counter-intuitive, just like in the explanation provided by annotators in Fig \ref{figure_main}.
% \\At the same time, for three different categories of arts, we designed different subtasks. For humorous videos, we set up a subtask to generate funny video titles, inspired by those interesting titles on social media. For Kasou Taishou videos, there are also subtasks that generate caption. In addition, in Kasou Taishou, each performance will have an official rating, and we use this rating as the fifth subtask of creative video: rating the creativity of the video. For magic video, we hope that the model can understand the principle of magic, so the subtask is to explain the realization of magic effect.\\
% Unlike most Video QA datasets' multi-choice and open-ended questions [], our task requires the model to generate a free-text answer without templates or word bags.\\
% Our experiments show that after finetuning on this dataset , the performance of the model on the unknown counter-intuitive video has improved, which means the model can "understand" the jokes in the video better.\\
% (Limitations) However, some models failed on xxx tasks...\\
% We also show the performance of the model on other tasks after FunQA adjustment.\\
% Finally, we hope that our work can become a practical and challenging benchmark, which explores the boundaries of video reasoning with its rich content, and inspires future work. Our benchmark can be found at xxxxxx[link].\\

\section{Introduction}
The charm of surprising videos, be they funny, creative, or filled with visual illusions, % always --- I guess videos have not always existed :D 
% been reflected
%a rich source of entertainment widely enjoyed by humanity.
% These videos possess a universal allure,
offer enjoyment and commands attention from viewers.
This type of media elicits \emph{positive surprise} \cite{noordewier2013valence},\footnote{c.f., \emph{negative} surprise, e.g., a surprising medical bill.}
a captivating emotion that stems not merely from perceiving surface-level visual stimuli, but rather, the innate ability of humans to understand and find delight in unexpected and counter-intuitive moments \cite{martin1983humour}. However, despite significant advancements in today's computer vision models, the question remains: can video models ``understand" the humor/creativity in surprising videos? %, decipher the humor, uncover the creativity, and appreciate the magic in the same way that we do?

Consider the humorous video depicted in Figure~\ref{figure_main} (left) as an illustrative example. We witness a man engrossed in his phone, sharing a meal with friends. Suddenly, one of his companions squeezes a generous amount of ketchup, which, instead of adorning the fries, splatters onto the man's face. The shock in his eyes, combined with his ketchup-covered visage, elicits laughter.\footnote{The hostility/superiority theory of humor posits that
humor can arise from claiming superiority over someone or something \cite{gruner1978understanding,billig2005laughter}; but alternate (more optimistic) theories of humor exist, \cite{attardo2008primer} offers a survey.} While humans effortlessly recognize this as an unusual (and potentially entertaining) event, % it poses a considerable challenge for computer vision to grasp its underlying humor.
% While models At the very least,
the reasoning required to holistically understand the scene is complex:
a model needs to % comprehend the context of the video, where
recognize that individuals were \textit{gathered to enjoy a meal together}, and discern that the comedic element arises from the \textit{ketchup intended for the fries ending up on the man's face instead}, and that the innocent expression of the unfortunate victim indicates no significant harm was caused.

While there have been some efforts to enhance computer vision models' performance in Video Question Answering (VideoQA), these works have primarily focused on the common, less surprising videos found in existing VideoQA datasets. Examples of commonly employed VideoQA datasets include YouCook2 \cite{zhou2018towards} which contains video clips from 2K cooking videos, Howto100M \cite{miech2019howto100m} which consists of only instructional videos.
While there exist video datasets that explore the humor in TV shows \cite{mustard, hasan-etal-2019-ur} and include tasks such as predicting laughter tracks \cite{Patro_2021_WACV}, these tasks often heavily rely on audio and narrative cues, with visual clues might playing a lesser role.
% TVQA \cite{lei2018tvqa} and other datasets \cite that cover TV shows may contain surprising scenarios, but the questions either mostly focused on literal factors, e.g., who is talking to who. 
% \jack{is this true?}
% These benchmarks % fall short when it comes to assessing a。 deeper level of human understanding,
% may partly % particularly in comprehending videos that are intriguing or captivating.
% As a result,
% Models trained on these datasets may not access ``unusual" videos, resulting in limited video reasoning capability.
%For example, datasets like TVQA \cite{lei2018tvqa} include factual questions based on TV episodes, while NExT-QA \cite{xiao2021nextqanext} involves describing and explaining temporal actions. 

To address this gap and evaluate computer vision models' ability to identify and understand visual commonsense violations in videos, we introduce \textbf{FunQA}, a comprehensive and high-quality VideoQA dataset comprising 4.3K surprising videos and 312K manually annotated free-text QA pairs. Our dataset consists of three subsets: \emph{1) HumorQA}, \emph{2) CreativeQA}, and \emph{3) MagicQA}. Each subset covers different sources and video contents, but the commonality lies in their surprising nature, e.g., the unexpected contrasts in humorous videos, the intriguing disguises in creative videos, and the seemingly impossible performances in magic videos. Our experiments suggest that these surprising videos require different types of reasoning than common videos, as existing VideoQA methods perform poorly on the corpus. With FunQA, we hope to provide a benchmark that covers the popular, important, and sophisticated genre of counter-intuitive/surprising videos. %, but also, is offers a fun challenge for existing models. %exudes a fascinating and pleasing charm.

%\jingkang{introduce 3 video types and why we include them}


In FunQA, we formulate three rigorous tasks to measure models' understanding of surprise:
\textit{1) Counter-intuitive timestamp localization}: a model must % This task involves
identify the specific time period within a video when an unexpected event takes place.
\textit{2) Detailed video description}: a model must generate coherent and objective descriptions of the video content, evaluating models' fundamental video understanding capabilities.
\textit{3) Counter-intuitiveness reasoning}: a model must generate concrete explanations of why the video is surprising. % which requires deep reasoning about the counter-intuitive events in the video.
These tasks progressively assess the model's ability to perceive, articulate, and reason about the counter-intuitive elements present in surprising videos. Additionally, we propose auxiliary tasks that pose higher-level challenges within the benchmark including assigning an appropriate and vivid title to the video, etc. 
%These challenges include assigning an appropriate and vivid title to the video and evaluating its creativity, where the model assesses the level of ingenuity and novelty displayed in the video content.
% Contrary to traditional video QA datasets that depend on multiple-choice and open-ended questions, our tasks require the model to generate free-text responses. We've been cautious in this journey to avoid content that could be counter-intuitive yet distasteful. The annotations we collect are of high quality and objectivity by employing cross-validation. For each art, we set up a subset to balance the different tastes and agreements people have with counter-intuitiveness.
To summarize our contributions: 

\textbf{1) New VideoQA Dataset:} We build a large-scale dataset \textbf{FunQA}, which complements the existing VideoQA dataset with intriguing videos.

\textbf{2) Novel and Challenging Tasks:} We design a number of novel tasks that allow the model to explore previously untouched problems, such as timestamp localization, and reasoning around counter-intuitiveness. These tasks push video reasoning beyond superficial descriptions, demanding deeper understanding and discernment.

\textbf{3) Comprehensive Evaluation:} We have done an extensive and comprehensive evaluation of cutting-edge baselines, giving the field an insight and future research direction. 

% \begin{itemize}
%     \item \textbf{New VideoQA Dataset:} We build a large-scale dataset \textbf{FunQA}, which complements the existing VideoQA dataset with intriguing videos.
%     \item \textbf{Novel and Challenging Tasks:} We design a number of novel tasks that allow the model to explore previously untouched problems, such as timestamp localization, and reasoning around counter-intuitiveness. These tasks push video reasoning beyond superficial descriptions, demanding deeper understanding and discernment. 
%     \item \textbf{Comprehensive Evaluation:} We have done an extensive and comprehensive evaluation of cutting-edge baselines, giving the field an insight and future research direction. 
% \end{itemize}
%Although FunQA enhances the model's performance on counter-intuitive videos, the journey to fully comprehending surprising content is still far from over. Significant performance gaps exist between model and human responses, showing models' limitations in spatial-temporal reasoning, visual-centered reasoning, and free-text generation. With FunQA, we aspire to provide a benchmark that pushes the limits of video reasoning, opens new avenues of exploration, and fuels progress in AI's capability to understand surprising videos. We believe that FunQA is a step forward towards endowing AI with a more human-like comprehension of the world. Our FunQA dataset and related resources are available at \href{https://github.com/Jingkang50/FunQA}{https://github.com/Jingkang50/FunQA}.

% In total, the FunQA benchmark is a challenging and high-quality VideoQA dataset specifically designed to evaluate and enhance the depth of video reasoning based on counter-intuitive and fun videos containing 312K free-text QA pairs spread across 24 video hours.\\
% We find that, while fine-tuning on FunQA does improve the model's performance on counter-intuitive videos, the journey to fully comprehending surprising content is a long one. There are significant performance gaps, especially in areas of spatial-temporal reasoning, visual-centered reasoning, and free-text generation.With FunQA, we hope to provide a benchmark that tests the boundaries of video reasoning, opens up new avenues of exploration, and sparks progress in AI's capability of understanding surprising videos. We believe that FunQA is a stepping stone towards making AI not just smarter, but also more human-like in its understanding of the world. Our FunQA dataset and related resources are available at \href{https://github.com/Jingkang50/FunQA}{https://github.com/Jingkang50/FunQA}. 
%Counter-intuitive experiences, those which subvert the norm, evoke a spectrum of emotions from curiosity and surprise to amusement or even annoyance. We've been cautious in this journey to avoid content that could be counter-intuitive yet distasteful. We guarantee that each video, despite being counter-intuitive, exudes a fascinating and pleasing charm. This aligns perfectly with our fundamental objective: for the AI model to comprehend these surprising videos and experience the same level of enjoyment as humans do.\\